% ============================================================
%  ArcLoom Processor Architecture v4.0
%  Pure Ternary Invariant Fabric & Hybrid Control Plane
%  ISO-Grade Deterministic Specification
%  Overleaf-compatible LaTeX
% ============================================================
\documentclass[11pt, a4paper]{article}

\usepackage[margin=2.2cm]{geometry}
\usepackage{amsmath, amssymb, amsthm, mathtools}
\usepackage{booktabs, array, longtable, tabularx, multirow}
\usepackage{xcolor}
\usepackage{hyperref}
\usepackage{titlesec}
\usepackage{enumitem}
\usepackage{fancyhdr}
\usepackage{caption}
\usepackage{float}

\hypersetup{
  colorlinks=true,
  linkcolor=black,
  urlcolor=blue,
}

\definecolor{al-gold}{RGB}{180, 140, 60}
\definecolor{al-dark}{RGB}{30, 30, 30}
\definecolor{al-green}{RGB}{40, 120, 60}
\definecolor{al-red}{RGB}{160, 40, 40}
\definecolor{al-gray}{RGB}{100, 100, 100}
\definecolor{al-blue}{RGB}{40, 80, 160}

\pagestyle{fancy}
\fancyhf{}
\rhead{\textcolor{al-gray}{\small ArcLoom Processor Specification v4.0}}
\lhead{\textcolor{al-gray}{\small Confidential --- Internal Use Only}}
\rfoot{\textcolor{al-gray}{\small Page \thepage}}

\titleformat{\section}{\large\bfseries\color{al-dark}}{}{0em}{}[\titlerule]
\titleformat{\subsection}{\normalsize\bfseries\color{al-dark}}{}{0em}{}
\titleformat{\subsubsection}{\normalsize\itshape\color{al-dark}}{}{0em}{}

\setlength{\parindent}{0pt}
\setlength{\parskip}{6pt}

\newtheorem{definition}{Definition}[section]
\newtheorem{axiom}{Axiom}[section]
\newtheorem{theorem}{Theorem}[section]
\newtheorem{corollary}{Corollary}[section]
\newtheorem{requirement}{Requirement}[section]

\newcolumntype{Y}{>{\raggedright\arraybackslash}X}

% ============================================================
\begin{document}

\begin{center}
  {\LARGE \textbf{ArcLoom Processor Architecture}}\\[6pt]
  {\large Pure Ternary Invariant Fabric \& Hybrid Control Plane}\\[4pt]
  {\large ISO-Grade Deterministic Specification}\\[4pt]
  {\large Version 4.0}\\[8pt]
  {\color{al-gray}\small April 2026 \quad|\quad Confidential --- Internal Use Only}\\[4pt]
  {\color{al-gray}\small Joseph Forrester \quad|\quad ArcLoom Systems / DSF-AI Architecture}
\end{center}

\vspace{1em}

\renewcommand{\arraystretch}{1.15}
\begin{tabularx}{\textwidth}{|p{0.28\textwidth}|Y|}
\hline
Document ID & ARC-PROC-SPEC-4.0 \\\hline
Version & 4.0 \\\hline
Status & Controlled specification \\\hline
Effective Date & 2026-04-10 \\\hline
Classification & Confidential --- Internal Use Only \\\hline
Purpose & Define the ArcLoom processor architecture in both Pure Ternary
          and Hybrid configurations, including physical geometry, deterministic
          logic, error hardening, L6-TCL integration, and reference platform
          design for autonomous robotic validation. \\\hline
Primary Audience & Hardware architects, FPGA engineers, embedded systems
                   engineers, DSF-AI integration engineers \\\hline
Change Control & Semantic changes require revision-history entry and
                 impact assessment \\\hline
\end{tabularx}

\vspace{1em}
\hrule
\vspace{1em}

\tableofcontents
\clearpage

% ============================================================
% PART I: FOUNDATIONS
% ============================================================

\section{Introduction}

\subsection{Purpose}

The ArcLoom processor is a non-von Neumann computing architecture designed
to serve as the hardware substrate for DSF-AI (Deterministic Structural
Field AI). It replaces binary transistor switching with topological phase
invariants, enabling direct physical computation of structural state
without instruction-fetch-decode-execute overhead.

This specification defines two configurations:
\begin{enumerate}[label=\arabic*.]
  \item \textbf{Configuration A: Pure Ternary Invariant Fabric} ---
        all computation performed through topological phase coupling.
        No binary logic. No clock. No instruction set. The processor
        \textit{is} the structural field.
  \item \textbf{Configuration B: Hybrid ArcLoom Processor Platform (AHPP)} ---
        a ternary invariant fabric supervised by a binary control coprocessor
        (BCP) for industrial-grade error correction and validation.
\end{enumerate}

\subsection{Scope}

This document covers:
\begin{itemize}[nosep]
  \item Mathematical foundations (balanced ternary, winding numbers, phase coupling)
  \item Physical geometry (triple-node resonant ring)
  \item Deterministic logic (resonance junctions, coupling matrix)
  \item Error hardening (metastability resolution, parity chains, H1-C policy)
  \item L6-TCL integration (dimensional grinder, structural lock)
  \item I/O boundary (krimelack transduction rings)
  \item Reference platform design (autonomous wheeled robot)
  \item Bill of materials and build specifications
\end{itemize}

\subsection{Design Constraints (Diamond Hard)}

\begin{enumerate}[label=\arabic*.]
  \item \textbf{Zero heuristics.} All thresholds derived from geometry or
        frozen rational constants.
  \item \textbf{Zero randomness.} No stochastic elements in the compute fabric.
  \item \textbf{Deterministic.} Given identical inputs, the processor produces
        identical outputs regardless of manufacturing variation (within
        tolerance bounds).
  \item \textbf{Domain agnostic.} The fabric computes structural state from
        any temporal signal; the domain is encoded in the coupling matrix,
        not the hardware.
\end{enumerate}

\subsection{Theoretical Basis}

The ArcLoom architecture rests on three established mathematical foundations:

\begin{enumerate}[label=\arabic*.]
  \item \textbf{Topological invariance:} Winding numbers of continuous phase
        fields are integer-valued and resistant to continuous deformation
        (noise). A winding number cannot change without a topological
        discontinuity (``tearing'') of the phase field.
  \item \textbf{Balanced ternary arithmetic:} The digit set $\{-1, 0, +1\}$
        provides symmetric signed representation, efficient carry propagation,
        and a native zero state (quiescence).
  \item \textbf{Coupled oscillator dynamics:} Networks of phase-coupled
        oscillators naturally minimize energy, producing deterministic
        steady-state configurations that encode computational results.
\end{enumerate}

% ============================================================
\section{Glossary of Terms}

\begin{longtable}{@{}p{4cm}p{10.5cm}@{}}
\toprule
\textbf{Term} & \textbf{Definition} \\
\midrule
\endhead
AIF & ArcLoom Invariant Fabric. The ternary computing substrate. \\
AHPP & ArcLoom Hybrid Processor Platform. AIF + BCP combined system. \\
ATF & ArcLoom Topology Fabric (synonym for AIF). \\
Balanced Ternary & Number system with digits $\{-1, 0, +1\}$ (base 3). \\
BCP & Binary Control Coprocessor. Conventional MCU/FPGA supervisor in
      hybrid configuration. \\
Coupling Matrix ($J_{ij}$) & Matrix defining phase-coupling strengths between
                             oscillator rings. Encodes the computational
                             topology. \\
Dimensional Grinder & The L6-TCL mechanism that maps structural constraints
                      to dimensional reduction. \\
Doorway Event & A local collapse of phase amplitude below threshold $\epsilon$,
                indicating topological vulnerability. \\
H1-C & Hardening Level 1 with Correction. The deterministic error recovery
       policy in hybrid mode. \\
Krimelack & Sensory transduction ring. Converts analog environmental signals
            into trit-compatible phase modulation. \\
Metastability & The condition where the phase field sits at a saddle point
                between two valid trit states, producing no valid output. \\
Negative Space & Quiescent structural region that acts as an active potential
                 well, driving inevitable state transitions. \\
Parity Chain & A group of trits whose winding numbers must sum to a known
               topological constant, providing geometric error correction. \\
Resonance Junction & The ternary equivalent of a logic gate; computation
                     occurs through energy minimization of coupled rings. \\
SL-1 & Structural Lock Level 1. The L6-TCL output certifying geometric
       inevitability. \\
Topological Schmitt Trigger & Nonlinear coupling term that creates dead zones
                              around saddle points, forcing metastability
                              resolution. \\
Trit & Ternary digit. The fundamental unit of ArcLoom computation.
      $t \in \{-1, 0, +1\}$. \\
Winding Number ($w$) & Integer topological invariant of a phase field around
                       a ring. Encodes the trit state. \\
\bottomrule
\end{longtable}

% ============================================================
\section{Mathematical Foundations}

\subsection{Balanced Ternary Representation}

The fundamental unit of computation is the balanced ternary digit (trit):
\begin{equation}
  t \in \{-1, 0, +1\}
\end{equation}

An $n$-trit word $T = (t_0, t_1, \ldots, t_{n-1})$ has value:
\begin{equation}
  \text{val}(T) = \sum_{i=0}^{n-1} t_i \cdot 3^i
  \label{eq:trit-value}
\end{equation}
where $t_0$ is the least-significant trit.

\textbf{Properties:}
\begin{itemize}[nosep]
  \item Information density: $\log_2 3 \approx 1.585$ bits per trit
  \item Symmetric signed representation without a sign bit
  \item Negation by digit-wise sign flip: $\overline{t_i} = -t_i$
  \item Native zero state ($w = 0$) represents structural quiescence
\end{itemize}

\subsection{Winding Number (Continuous)}

\begin{definition}[Winding Number --- Continuous]
The winding number of a continuous phase field $\phi(\theta)$ around a closed
ring is:
\begin{equation}
  w = \frac{1}{2\pi} \oint \frac{d\phi}{d\theta}\, d\theta \in \mathbb{Z}
  \label{eq:winding-continuous}
\end{equation}
\end{definition}

The winding number is a \textbf{topological invariant}: it cannot change under
continuous deformation of the phase field. Only a discontinuity (``tear'')
in the field can alter $w$. This is the physical basis for noise resistance.

\subsection{Winding Number (Discrete)}

In a physical implementation with $N$ sample points around the ring:
\begin{equation}
  w = \frac{1}{2\pi} \sum_{i=0}^{N-1} \text{wrap}(\phi_{i+1} - \phi_i),
  \qquad \phi_N \equiv \phi_0
  \label{eq:winding-discrete}
\end{equation}
where $\text{wrap}(\Delta\phi)$ constrains the phase difference to $(-\pi, \pi]$.

\subsection{Phase Coupling Energy}

For a system of coupled oscillator rings, the total energy is:
\begin{equation}
  E_{sys} = -\sum_{i,j} J_{ij} \cos(\phi_i - \phi_j)
  \label{eq:coupling-energy}
\end{equation}
where $J_{ij}$ is the coupling strength between oscillators $i$ and $j$.
The system naturally evolves toward the energy minimum, producing
deterministic steady-state configurations.

\subsection{Nonlinear Coupling (Topological Schmitt Trigger)}

To resolve metastability at saddle points, a nonlinear term is added to
the coupling energy:
\begin{equation}
  E_{sys} = -\sum_{i,j} J_{ij} \left[ \cos(\Delta\phi_{ij})
            + \alpha \cos^3(\Delta\phi_{ij}) \right]
  \label{eq:schmitt}
\end{equation}
where:
\begin{itemize}[nosep]
  \item $\Delta\phi_{ij} = \phi_i - \phi_j$
  \item $\alpha$ is a frozen rational constant determining the dead-zone width
        around saddle points
\end{itemize}

The cubic term deepens the energy wells at valid trit states ($w \in \{-1,0,+1\}$)
and raises the barrier at saddle points, forcing rapid resolution of ambiguous
phase states.

\subsection{$n$-Sphere Volume and Gamma Function}

(Refer to L6-TCL Specification v1.0 for full derivation.)

\begin{equation}
  V_n(R) = \frac{\pi^{n/2}}{\Gamma\!\left(\frac{n}{2}+1\right)} R^n
\end{equation}

The capture threshold $\tau$ is derived from the gamma-function inflection:
\begin{equation}
  \tau = 1 - \frac{\Gamma\!\left(\frac{n_{start}}{2}+1\right)}
                  {\Gamma\!\left(\frac{n_{eff}}{2}+1\right)}
  \cdot \pi^{(n_{eff}-n_{start})/2}
\end{equation}

Structural Lock fires when $n_{eff} < n_{start}/e$.

% ============================================================
\section{Physical Geometry: The Triple-Node Resonant Ring}

\subsection{Description}

A single trit register consists of three nonlinear oscillators (Nodes $A$,
$B$, $C$) arranged in an equilateral triangle and coupled via phase-shifting
delay lines. The state of the register is the emergent winding number of
the phase field across the three nodes.

\subsection{State Encoding}

\begin{center}
\begin{tabular}{@{}clll@{}}
\toprule
\textbf{Winding} & \textbf{State} & \textbf{Phase Relationship} & \textbf{Physical Meaning} \\
\midrule
$w = 0$  & Quiescent   & $\phi_A = \phi_B = \phi_C$ & Uniform coherence; no flow \\
$w = +1$ & CW Flow     & $\phi_B = \phi_A + 2\pi/3$, $\phi_C = \phi_B + 2\pi/3$ & Clockwise phase progression \\
$w = -1$ & CCW Flow    & $\phi_C = \phi_A + 2\pi/3$, $\phi_B = \phi_C + 2\pi/3$ & Counter-clockwise progression \\
\bottomrule
\end{tabular}
\end{center}

\subsection{Stability Analysis}

Each winding state occupies a local energy minimum of the coupling potential.
Transitions between states require traversal of an energy barrier proportional
to $J_{ij}$ and amplified by the nonlinear Schmitt term ($\alpha$). Small
perturbations (noise) are insufficient to cross the barrier; only a structured
input (a genuine signal) can drive a state transition.

The energy barrier height between adjacent states is:
\begin{equation}
  \Delta E = J \cdot (1 + 3\alpha/4) \cdot (1 - \cos(2\pi/3))
           = J \cdot (1 + 3\alpha/4) \cdot 1.5
\end{equation}

For $J = 1$ and $\alpha = 37/64$: $\Delta E \approx 2.37$ (dimensionless
coupling units).

% ============================================================
\section{Configuration A: Pure Ternary Invariant Fabric}

\subsection{Description}

The Pure Ternary configuration removes all binary logic. Computation,
error correction, and I/O are performed entirely through topological
phase dynamics. There is no clock, no instruction set, and no
fetch-decode-execute cycle.

\subsection{Deterministic Logic via Resonance Junctions}

A \textbf{Resonance Junction} is the ternary equivalent of a logic gate.
Two or more input rings are coupled to one or more output rings. The
coupled system evolves to its energy minimum, producing deterministic
output states.

\subsubsection{Ternary Addition}

For a dual-ring adder computing $A + B = S + 3C$ (sum $S$, carry $C$):
\begin{itemize}[nosep]
  \item Input rings $R_A$ and $R_B$ hold winding numbers $w_A$ and $w_B$
  \item Output ring $R_S$ settles to $w_S = (w_A + w_B) \bmod 3_{\pm}$
        (balanced ternary modulo)
  \item Carry ring $R_C$ settles to $w_C = \lfloor(w_A + w_B + 1)/3\rfloor - 1$
\end{itemize}

The coupling matrix for this junction is:
\begin{equation}
  J = \begin{pmatrix}
    0 & 0 & J_{AS} & J_{AC} \\
    0 & 0 & J_{BS} & J_{BC} \\
    J_{AS} & J_{BS} & 0 & 0 \\
    J_{AC} & J_{BC} & 0 & 0
  \end{pmatrix}
\end{equation}

where coupling strengths are frozen rational constants derived from the
energy landscape of balanced ternary addition.

\subsection{The Coupling Matrix ($J_{ij}$): The Instruction Set}

In a von Neumann architecture, the instruction set defines what operations
the processor can perform. In ArcLoom, \textbf{the coupling matrix $J_{ij}$
is the instruction set}. It encodes:
\begin{itemize}[nosep]
  \item Which rings are coupled (topology)
  \item How strongly they are coupled (constraint strength)
  \item The direction of coupling (symmetric or asymmetric)
\end{itemize}

\begin{requirement}[Coupling Matrix Derivation]
$J_{ij}$ \textbf{must} be derived from the DSF-AI kernel structure. Each
L0--L4 output maps to a specific coupling strength between specific ring
pairs:
\begin{center}
\begin{tabular}{@{}lll@{}}
\toprule
\textbf{Kernel Output} & \textbf{Coupling Target} & \textbf{$J_{ij}$ Derivation} \\
\midrule
$D_k$ (Displacement) & Direction ring pair & $J = D_k \cdot J_{base}$ \\
$S_{UF}$ (Stability) & Convergence ring pair & $J = S_{UF} \cdot J_{base}$ \\
$M_k$ (Motion) & Momentum ring pair & $J = |M_k| \cdot J_{base}$ \\
$C_k$ (Cohesion) & Binding ring pair & $J = \frac{C_k}{1+C_k} \cdot J_{base}$ \\
$P_k$ (Pressure) & Compression ring pair & $J = \frac{P_k}{1+P_k} \cdot J_{base}$ \\
$B_k$ (Breathing) & Conviction ring pair & $J = |B_k| \cdot J_{base}$ \\
$U^*_k$ (Uncertainty) & Freedom ring pair & $J = -(U^*_k) \cdot J_{base}$ \\
$R_{rev,k}$ (Reversal) & Path-kill ring pair & $J = R_{rev,k} \cdot J_{max}$ \\
\bottomrule
\end{tabular}
\end{center}
$J_{base}$ is a frozen rational constant. $J_{max}$ is the maximum coupling
strength (forces immediate state transition).
\end{requirement}

\subsection{Geometric Error Correction: Topological Parity Chains}

\begin{definition}[Parity Chain]
A parity chain is a group of $P$ trits whose winding numbers must sum
to a known topological constant $K$:
\begin{equation}
  \sum_{i=1}^{P} w_i = K
  \label{eq:parity}
\end{equation}
where $K$ is fixed at fabrication time and is itself a topological invariant.
\end{definition}

If manufacturing drift or noise causes one trit to wander, the parity
constraint from its neighbors creates a restoring force:
\begin{equation}
  F_{restore,i} = -\lambda_{parity} \left( \sum_{j=1}^{P} w_j - K \right)
  \label{eq:parity-force}
\end{equation}
where $\lambda_{parity}$ is a frozen rational coupling constant.

\textbf{Recommended chain length:} $P = 5$. This provides single-trit
error correction capability (any single trit error is detectable and
correctable from the remaining 4).

\subsection{Metastability Resolution: Topological Schmitt Trigger}

When the phase field enters the saddle region between two valid states,
the nonlinear coupling term (Equation~\ref{eq:schmitt}) creates an
unstable equilibrium that forces resolution within a bounded time:

\begin{equation}
  t_{resolve} < \frac{1}{J \cdot \alpha \cdot \omega_0}
\end{equation}

where $\omega_0$ is the natural oscillator frequency. For typical parameters
($J = 1$, $\alpha = 37/64$, $\omega_0 = 2\pi \times 100\text{ MHz}$):
$t_{resolve} < 2.7$ ns.

\subsection{Power Consumption Model}

Power consumption is proportional to state-transition density, not clock
frequency:
\begin{equation}
  P_{avg} = N_{rings} \cdot \rho_{transition} \cdot E_{transition}
\end{equation}
where:
\begin{itemize}[nosep]
  \item $N_{rings}$ = total ring count
  \item $\rho_{transition}$ = state transitions per second (signal-dependent)
  \item $E_{transition}$ = energy per state transition (coupling-dependent)
\end{itemize}

During quiescence ($\rho_{transition} \to 0$), power consumption approaches
the static leakage floor. For structural detection tasks where most time
is spent in quiescence, this is inherently efficient.

% ============================================================
\section{Configuration B: Hybrid ArcLoom Processor Platform (AHPP)}

\subsection{Description}

The AHPP pairs the ArcLoom Invariant Fabric (AIF) with a Binary Control
Coprocessor (BCP). The BCP supervises timing, error correction, memory
management, and safety monitoring. The AIF performs structural computation.

This configuration is recommended for initial validation, commercial
deployment, and domains requiring certified determinism (medical, automotive).

\subsection{Architecture}

\begin{center}
\begin{tabular}{@{}lll@{}}
\toprule
\textbf{Component} & \textbf{Function} & \textbf{Implementation} \\
\midrule
AIF (Invariant Fabric) & Structural computation & FPGA phase field (Option P1) \\
                       &                        & or analog oscillator array (Option P2) \\
BCP (Binary Coprocessor) & Supervision, correction & MCU (STM32H7 or equivalent) \\
Structural Timebase & Monotonic timing & Crystal oscillator + counter \\
Doorway Detector & Amplitude monitoring & Comparator array \\
\bottomrule
\end{tabular}
\end{center}

\subsection{Doorway Detection (Negative Space Monitoring)}

The BCP continuously monitors the amplitude $a_i$ of each oscillator node.
A doorway event occurs when amplitude collapses below threshold:

\begin{equation}
  \delta_j^{ext} = \begin{cases}
    1 & \text{if } \min(a_i) < \epsilon_{gate} \\
    0 & \text{otherwise}
  \end{cases}
\end{equation}

where $\epsilon_{gate}$ is a frozen threshold derived from the coupling
energy landscape.

\subsection{H1-C Hardening and Correction Policy}

Upon detecting a doorway event, the BCP executes deterministic recovery:

\begin{enumerate}[label=\arabic*.]
  \item \textbf{STALL:} Computation halts. Event logged with structural
        timestamp.
  \item \textbf{Amplitude Healing:} The BCP restores amplitude at the slip
        site to baseline.
  \item \textbf{Relaxation:} $R$ extra relaxation cycles are executed
        ($R$ is a frozen integer constant).
  \item \textbf{Verification:} Winding number is re-read.
        \begin{itemize}[nosep]
          \item If $w$ matches pre-stall value: \textbf{Recovered STALL}.
          \item If $w$ differs but parity chain is consistent: \textbf{Corrected}.
          \item If parity chain is inconsistent: \textbf{Binary Fallback} ---
                the BCP recomputes the kernel in binary to ensure output
                determinism.
        \end{itemize}
\end{enumerate}

\subsection{Binary Fallback}

The BCP maintains a complete binary implementation of the L0--L4 kernel.
If the AIF produces an uncorrectable result, the BCP substitutes its own
binary-computed output. This guarantees deterministic output under all
conditions, at the cost of losing the efficiency advantage for that
computation cycle.

% ============================================================
\section{I/O Boundary: Krimelack Transduction Rings}

\subsection{Description}

The krimelack is the sensory interface between the physical environment
and the trit fabric. It converts analog environmental signals into
trit-compatible phase modulation without requiring analog-to-digital
conversion.

\subsection{Operating Principle}

A krimelack is an oscillator ring whose natural frequency is modulated
by an external physical signal (light intensity, temperature, pressure,
distance, etc.). The modulated frequency shifts the ring's phase
relationship with the fabric, changing its effective winding number.

\begin{equation}
  \omega_{krim}(t) = \omega_0 + \kappa \cdot s(t)
\end{equation}

where:
\begin{itemize}[nosep]
  \item $\omega_0$ = natural frequency (matched to fabric)
  \item $\kappa$ = transduction coefficient (signal-type dependent)
  \item $s(t)$ = raw analog sensor signal
\end{itemize}

The phase accumulation over time naturally produces winding-number
transitions when the signal crosses structural thresholds:

\begin{equation}
  \Delta w = \left\lfloor \frac{1}{2\pi} \int_0^T [\omega_{krim}(t) - \omega_0]\, dt + \frac{1}{2} \right\rfloor
\end{equation}

\subsection{Sensor Types and Transduction}

\begin{center}
\begin{tabular}{@{}llll@{}}
\toprule
\textbf{Sensor} & \textbf{Signal $s(t)$} & \textbf{Transduction} & \textbf{Krimelack Type} \\
\midrule
Vision (camera) & Pixel intensity & Photoresistor $\to \omega$ modulation & Optical krimelack \\
Depth (ToF/LIDAR) & Distance & Capacitive $\to \omega$ modulation & Range krimelack \\
Motion (IMU) & Acceleration & Piezo $\to \omega$ modulation & Inertial krimelack \\
Proximity (IR/US) & Object distance & Reflective $\to \omega$ modulation & Proximity krimelack \\
Temperature & Thermal & Thermistor $\to \omega$ modulation & Thermal krimelack \\
\bottomrule
\end{tabular}
\end{center}

\subsection{Output Transduction}

The output boundary converts trit states into physical actuator commands.
Motor control is achieved by coupling output rings to PWM generators
(hybrid) or directly to motor driver oscillators (pure):

\begin{equation}
  v_{motor} = w_{out} \cdot v_{max} \qquad (w_{out} \in \{-1, 0, +1\})
\end{equation}

This produces: reverse ($-1$), stop ($0$), forward ($+1$). Proportional
control is achieved through trit-word encoding (multiple trits per motor
channel).

% ============================================================
\section{L6-TCL Integration: The Dimensional Grinder}

\subsection{Description}

The ArcLoom processor physically instantiates the L6 Topological Constraint
Layer. The fabric's internal state space \textit{is} the $n$-dimensional
manifold. As constraints are applied through the coupling matrix, the
effective dimensionality of the fabric's state space reduces.

\subsection{Constraint Mapping}

Each L0--L4 kernel output modifies the coupling matrix, which in turn
restricts the fabric's available phase configurations:

\begin{equation}
  n_{eff} = n_{start} - \sum_{i=1}^{k} \text{rank}(C_i)
\end{equation}

where $n_{start}$ is the total ring count (the fabric's maximum
dimensionality) and each constraint $C_i$ removes one or more degrees
of freedom by locking specific ring-pair phase relationships.

\subsection{Structural Lock (SL-1)}

When $n_{eff} < n_{start}/e$, the fabric physically enters the capture
basin. The remaining phase configurations cannot sustain an alternative
output. The processor issues SL-1: a hardware interrupt indicating
geometric certainty.

In the pure configuration, SL-1 manifests as a \textbf{global phase lock}
--- all rings simultaneously settle to a single coherent state. In the
hybrid configuration, the BCP detects this condition and asserts a
hardware interrupt line.

% ============================================================
\section{Reference Platform: Autonomous Wheeled Robot}

\subsection{Purpose}

The reference platform is a wheeled robot that demonstrates ArcLoom's
structural perception capabilities in a physical environment. The robot
can:
\begin{enumerate}[nosep]
  \item Follow a human operator (person-following mode)
  \item Navigate autonomously around obstacles (exploration mode)
  \item Follow a predefined course (course-following mode)
\end{enumerate}

This validates the L0--L4 kernel operating on real sensory data through
real krimelack transduction, processed by a real ArcLoom fabric, producing
real motor commands.

\subsection{Platform Specifications}

\begin{center}
\begin{tabular}{@{}lll@{}}
\toprule
\textbf{Parameter} & \textbf{Specification} & \textbf{Notes} \\
\midrule
Chassis & 4WD wheeled platform, $\sim$30cm $\times$ 25cm & Aluminum or 3D-printed \\
Drive & Differential drive, 2 powered + 2 caster & DC gear motors with encoders \\
Speed & 0.5--1.5 m/s & Configurable via trit-word \\
Battery & LiPo 3S (11.1V), 5000 mAh & $\sim$2 hours runtime \\
Weight & $<$ 3 kg (including processor and sensors) & \\
\bottomrule
\end{tabular}
\end{center}

% ============================================================
\subsection{Configuration A Build: Pure Ternary Robot}

\subsubsection{Processor}

\begin{center}
\begin{tabular}{@{}llll@{}}
\toprule
\textbf{Component} & \textbf{Part} & \textbf{Function} & \textbf{Qty} \\
\midrule
FPGA (Phase Fabric) & Xilinx Artix-7 (XC7A200T) & Trit fabric simulation & 1 \\
                    &                             & 200K logic cells, 740 DSP & \\
Oscillator Array & Si5351A PLL clock gen & Phase reference generation & 4 \\
DAC (Phase Output) & AD5686R quad 16-bit DAC & Phase-to-analog for motors & 2 \\
ADC (Krimelack Input) & AD7689 8-ch 16-bit ADC & Sensor-to-phase conversion & 2 \\
Power Regulator & TPS65281 multi-rail & 1.0V core, 1.8V I/O, 3.3V sensor & 1 \\
\bottomrule
\end{tabular}
\end{center}

\subsubsection{Sensor Array (Krimelacks)}

\begin{center}
\begin{tabular}{@{}lllll@{}}
\toprule
\textbf{Sensor} & \textbf{Part} & \textbf{Interface} & \textbf{Range} & \textbf{Qty} \\
\midrule
Stereo Vision & OV7675 dual camera module & Parallel/SPI & 640$\times$480, 30fps & 1 pair \\
Depth/ToF & VL53L5CX (8$\times$8 zone) & I$^2$C & 2--400cm & 2 (front+rear) \\
IMU & ICM-42688-P 6-axis & SPI & $\pm$16g, $\pm$2000 dps & 1 \\
IR Proximity & VCNL4040 & I$^2$C & 0--20cm & 4 (corners) \\
Wheel Encoder & AS5048A magnetic & SPI & 14-bit absolute & 2 \\
\bottomrule
\end{tabular}
\end{center}

\subsubsection{Actuators}

\begin{center}
\begin{tabular}{@{}llll@{}}
\toprule
\textbf{Component} & \textbf{Part} & \textbf{Specification} & \textbf{Qty} \\
\midrule
Drive Motor & Pololu 37D 12V gear motor & 200 RPM, encoder & 2 \\
Motor Driver & DRV8874 H-Bridge & 6A peak, PWM input & 2 \\
Servo (camera pan) & SG90 micro servo & 180$^\circ$, 1.5kg-cm & 1 \\
\bottomrule
\end{tabular}
\end{center}

\subsubsection{Programming the Pure Ternary Fabric}

The pure ternary ArcLoom is not ``programmed'' in the conventional sense.
It is \textbf{configured} by defining:

\begin{enumerate}[label=\arabic*.]
  \item \textbf{Ring Topology:} The number and arrangement of trit rings
        in the FPGA. For the reference robot: 64 trit rings organized in
        an 8$\times$8 lattice (matching the kernel's 8-dimensional state
        vector, replicated across 8 spatial zones).
  \item \textbf{Coupling Matrix ($J_{ij}$):} Loaded into the FPGA as a
        lookup table of fixed-point coupling coefficients. Derived from
        the L0--L4 kernel structure per the mapping in Section~5.3.
  \item \textbf{Krimelack Binding:} Each sensor ADC channel is bound to
        a specific input ring. The transduction coefficient $\kappa$ is
        set per sensor type.
  \item \textbf{Parity Chain Assignment:} Every group of 5 adjacent rings
        is assigned a parity constant $K$.
  \item \textbf{Output Binding:} Specific output rings are coupled to
        motor driver DAC channels.
\end{enumerate}

\textbf{There is no instruction memory, no program counter, no opcodes.}
The FPGA implements the phase-coupling physics as a dataflow fabric. Input
signals propagate through the coupling network and produce output states
through energy minimization.

The ``behavior'' of the robot is entirely determined by the coupling
matrix. Different coupling matrices produce different behaviors (follow,
explore, course-follow) without changing the hardware.

% ============================================================
\subsection{Configuration B Build: Hybrid Robot}

\subsubsection{Processor}

\begin{center}
\begin{tabular}{@{}llll@{}}
\toprule
\textbf{Component} & \textbf{Part} & \textbf{Function} & \textbf{Qty} \\
\midrule
BCP (Binary Supervisor) & STM32H743 MCU & Supervision, H1-C, fallback & 1 \\
                        & 480 MHz Cortex-M7 & Binary kernel backup & \\
AIF (Invariant Fabric) & Xilinx Artix-7 (XC7A100T) & Phase fabric (smaller) & 1 \\
                       & 100K logic cells & \\
Clock Generator & Si5351A & Structural timebase & 1 \\
ADC & ADS8688 8-ch 16-bit & Sensor input & 1 \\
DAC & DAC8568 8-ch 16-bit & Motor/servo output & 1 \\
\bottomrule
\end{tabular}
\end{center}

\subsubsection{Sensor Array}

Same sensor complement as Configuration A (Section 10.2.2). The BCP
handles sensor initialization, calibration, and I$^2$C/SPI bus management.

\subsubsection{Programming the Hybrid Processor}

The hybrid ArcLoom is programmed in two layers:

\textbf{Layer 1: AIF Configuration (same as pure)}
\begin{itemize}[nosep]
  \item Ring topology, coupling matrix, krimelack binding, parity chains,
        output binding --- all loaded into the FPGA as fixed-point coefficients.
\end{itemize}

\textbf{Layer 2: BCP Firmware (C/Rust)}
\begin{itemize}[nosep]
  \item \texttt{main\_loop()}: Read AIF state $\to$ doorway detection $\to$
        H1-C correction $\to$ write outputs
  \item \texttt{l0\_l4\_kernel()}: Binary backup implementation of the full
        DSF-AI kernel (used only during binary fallback)
  \item \texttt{sensor\_init()}: I$^2$C/SPI bus setup, calibration sequences
  \item \texttt{safety\_monitor()}: Watchdog, battery voltage, motor current limits
  \item \texttt{telemetry()}: UART/BLE output for debugging and monitoring
\end{itemize}

The BCP firmware is deterministic: no dynamic memory allocation, no
floating-point (fixed-point only), no interrupts during critical sections,
no RTOS (bare-metal superloop).

% ============================================================
\section{Comparison of Configurations}

\begin{center}
\begin{tabular}{@{}p{3.5cm}p{5.5cm}p{5.5cm}@{}}
\toprule
\textbf{Feature} & \textbf{Config A: Pure Ternary} & \textbf{Config B: Hybrid AHPP} \\
\midrule
Logic basis & Topological invariants only & Binary-supervised phases \\
Error correction & Geometric (parity chains + Schmitt) & H1-C + binary fallback \\
Clock & Asynchronous (structural) & BCP clock + structural fabric \\
Power (quiescent) & Near-zero (static leakage only) & BCP idle + fabric leakage \\
Power (active) & Proportional to transitions & BCP active + fabric transitions \\
Determinism guarantee & Topological (geometric) & BCP-certified (computational) \\
Certification path & Novel (no precedent) & Standard (MCU-based safety case) \\
Manufacturing tolerance & Parity chains required & BCP compensates drift \\
Debugging & Field observation only & Register inspection + telemetry \\
Recommended use & Research, sentience exploration & Commercial, medical, automotive \\
Development timeline & 12--18 months & 6--9 months \\
\bottomrule
\end{tabular}
\end{center}

% ============================================================
\section{Recommended Build Sequence}

\begin{enumerate}[label=\textbf{Phase \arabic*:}]
  \item \textbf{FPGA Simulation (Months 1--3):} Implement the phase fabric
        in Verilog/VHDL on the Artix-7 using fixed-point arithmetic. Validate
        trit stability, parity chains, and Schmitt trigger behavior in
        simulation. No physical sensors.
  \item \textbf{Hybrid Robot Build (Months 3--6):} Integrate STM32H7 BCP
        with FPGA fabric. Connect sensors and motors. Validate H1-C policy
        with real noise. Demonstrate person-following and obstacle avoidance.
  \item \textbf{Pure Ternary Validation (Months 6--12):} Disable the BCP.
        Run the fabric standalone with parity chain error correction only.
        Characterize noise tolerance, metastability rate, and power consumption.
        Compare outputs to hybrid (binary fallback) baseline.
  \item \textbf{Analog Fabric (Months 12--18):} Replace FPGA simulation with
        physical oscillator arrays (PLL networks). This is the transition
        from ``simulated topology'' to ``real topology.'' Validate that
        physical phase dynamics match FPGA simulation.
\end{enumerate}

% ============================================================
\section{Invariants and Constraints}

\begin{enumerate}
  \item The winding number is the \textbf{only} valid representation of a
        trit state. Voltage levels, register values, or other binary
        representations are implementation artifacts, not architectural
        primitives.
  \item The coupling matrix $J_{ij}$ is \textbf{derived}, not fitted. Any
        implementation that tunes $J_{ij}$ based on output performance is
        non-compliant.
  \item Parity chain constants $K$ are set at fabrication/configuration
        time and are \textbf{never modified} during operation.
  \item The nonlinear coupling constant $\alpha$ (Schmitt trigger) is a
        frozen rational constant.
  \item In hybrid mode, the BCP \textbf{never modifies} the coupling matrix.
        It can only stall, heal amplitude, and substitute binary-computed
        output. It cannot alter the fabric's topology.
  \item SL-1 (Structural Lock) is a hardware output, not a software flag.
        Once asserted, it remains asserted until the next structural frame
        or exogenous disruption.
  \item The krimelack transduction coefficient $\kappa$ is a frozen constant
        per sensor type. It is not adapted at runtime.
\end{enumerate}

% ============================================================
\section{Version History}

\begin{center}
\begin{tabular}{@{}lll@{}}
\toprule
\textbf{Version} & \textbf{Date} & \textbf{Description} \\
\midrule
1.0 & 2026-02 & Initial ArcLoom concept (LoomCore/RNA3D). \\
2.0 & 2026-03 & Hybrid architecture with H1-C hardening. \\
3.0 & 2026-04 & Pure ternary model added (w3.1 collaboration). \\
4.0 & 2026-04-10 & Comprehensive ISO-grade specification. Added: coupling matrix \\
    &            & derivation, topological Schmitt trigger, parity chains, krimelack \\
    &            & I/O boundary, metastability resolution, reference robot platform \\
    &            & with BOM for both pure and hybrid configurations. \\
\bottomrule
\end{tabular}
\end{center}

\vspace{2em}
\hrule
\vspace{0.5em}
{\color{al-gray}\small This document is confidential and for internal use only.
It constitutes the ISO-grade deterministic specification for the ArcLoom
Processor Architecture in both Pure Ternary and Hybrid configurations.
All mathematical operations are deterministic. All constants are frozen
rationals or derived from natural mathematical constants ($e$, $\Gamma$,
$\pi$). No parameters are fitted, estimated, or heuristically determined.}

\end{document}
